%! Equivalent Representations of Trigonometric Functions — Unit 3, Lesson 3.12 — Group Activity
\documentclass[10pt]{article}
\usepackage{precalculus-article}
\usepackage{precalculus-boxes}
\newcommand{\TallMath}[1]{$\displaystyle #1\rule[-1.4em]{0pt}{3.2em}$}

\begin{document}

\pageheader{Unit 3, Lesson 3.12}{Group Activity}
\namepartnerperiod

\vspace{0.10in}

%----------------------------------------------------------------------
% Context
%----------------------------------------------------------------------
\begin{scenariobox}[Context: Solving Equations Using Trigonometric Identities]{sky}
Equations like $2\sin^2 x + \cos x = 1$ mix two different trig functions and
cannot be solved directly. In this activity your group will use the Pythagorean,
double-angle, and sum identities to rewrite such equations in terms of a
\emph{single} trig function, then solve them algebraically.
Choose \textbf{one tier} based on your readiness. Everyone should be prepared
to explain your strategy and verify your solutions.
\end{scenariobox}

\vspace{0.08in}

%----------------------------------------------------------------------
% Tier R
%----------------------------------------------------------------------
\begin{tcolorbox}[colback=white, colframe=black!40,
    title=\textbf{Tier R --- Reinforce},
    fonttitle=\bfseries, arc=2mm, left=3mm, right=3mm, top=2mm, bottom=2mm]

\textbf{Directions:} Solve $2\sin^2 x + \cos x - 1 = 0$ for $x \in [0, 2\pi)$.
The substitution $\sin^2 x = 1 - \cos^2 x$ is provided. Complete each step.

\vspace{6pt}
\textbf{Step 1.} Substitute $\sin^2 x = 1 - \cos^2 x$ into the equation:
\vspace{4pt}
$2(\;\underline{\hspace{3.5cm}}\;) + \cos x - 1 = 0$

\vspace{8pt}
\textbf{Step 2.} Expand:
$\underline{\hspace{2.5cm}} + \cos x - 1 = 0$

\vspace{8pt}
\textbf{Step 3.} Collect all terms on one side and multiply through by $-1$:
$\underline{\hspace{7cm}} = 0$

\vspace{8pt}
\textbf{Step 4.} Let $u = \cos x$. Rewrite as a quadratic in $u$:
$\underline{\hspace{3.5cm}} = 0$

\vspace{4pt}
Factor: $(\;\underline{\hspace{2cm}}\;)(\;\underline{\hspace{2cm}}\;) = 0$

\vspace{8pt}
\textbf{Step 5.} Solve each factor:
\vspace{4pt}

$\cos x = $ \underline{\hspace{2.5cm}} \qquad or \qquad $\cos x = $ \underline{\hspace{2.5cm}}

\vspace{8pt}
\textbf{Step 6.} Find all $x \in [0, 2\pi)$ for each case.

\vspace{4pt}
\renewcommand{\arraystretch}{1.8}
\begin{tabular}{|l|l|}
\hline
$\cos x = $ \underline{\hspace{1.8cm}} & Solutions: \underline{\hspace{3.5cm}} \\
\hline
$\cos x = $ \underline{\hspace{1.8cm}} & Solutions: \underline{\hspace{3.5cm}} \\
\hline
\end{tabular}

\vspace{8pt}
\textbf{Answer:} $x \in \big\{ \underline{\hspace{5.5cm}} \big\}$

\vspace{8pt}
\textbf{Check:} Substitute each solution back into the original equation to verify.

\vspace{4pt}
$x = 0$: $2\sin^2(0) + \cos(0) - 1 = $ \underline{\hspace{1.2cm}} $+$ \underline{\hspace{1.2cm}} $-$ 1 $=$ \underline{\hspace{1.2cm}}

\end{tcolorbox}

\vspace{0.08in}

%----------------------------------------------------------------------
% Tier A
%----------------------------------------------------------------------
\begin{tcolorbox}[colback=white, colframe=black!40,
    title=\textbf{Tier A --- Apply},
    fonttitle=\bfseries, arc=2mm, left=3mm, right=3mm, top=2mm, bottom=2mm]

\textbf{Directions:} Solve each equation for $x \in [0, 2\pi)$. Show all work.
State which identity you use in each problem.

\vspace{8pt}
\textbf{Problem 1.} $2\cos^2 x - \cos x - 1 = 0$

\textit{(No substitution needed --- already in one function. Factor directly.)}

\vspace{4pt}
Identity used: \underline{\hspace{4cm}}

\writelines{4}

\textbf{Solutions:} $x \in \big\{ \underline{\hspace{4cm}} \big\}$

\vspace{8pt}
\textbf{Problem 2.} $\sin 2x = \sin x$

\textit{(Hint: rewrite $\sin 2x$ using the double-angle identity.)}

\vspace{4pt}
Identity used: \underline{\hspace{4cm}}

\writelines{5}

\textbf{Solutions:} $x \in \big\{ \underline{\hspace{4cm}} \big\}$

\vspace{8pt}
\textbf{Problem 3.} $\cos^2 x - \sin^2 x = \dfrac{1}{2}$

\textit{(Hint: the left side equals a double-angle expression.)}

\vspace{4pt}
Identity used: \underline{\hspace{4cm}}

\writelines{5}

\textbf{Solutions:} $x \in \big\{ \underline{\hspace{4cm}} \big\}$

\end{tcolorbox}

\vspace{0.08in}

%----------------------------------------------------------------------
% Tier E
%----------------------------------------------------------------------
\begin{tcolorbox}[colback=white, colframe=black!40,
    title=\textbf{Tier E --- Extend},
    fonttitle=\bfseries, arc=2mm, left=3mm, right=3mm, top=2mm, bottom=2mm]

\textbf{Directions:} Complete each task. Show all algebraic steps.

\vspace{8pt}
\textbf{Part A --- Verify Identities.} Work from the \emph{left side only}; transform it
to match the right side.

\vspace{4pt}
\textbf{Identity 1:}\quad $\dfrac{\sin^2\theta}{1 - \cos\theta} = 1 + \cos\theta$

\vspace{4pt}
\textit{(Hint: substitute $\sin^2\theta = 1 - \cos^2\theta$, then factor.)}

\writelines{5}

\vspace{4pt}
\textbf{Identity 2:}\quad $\tan\theta + \cot\theta = \dfrac{1}{\sin\theta\cos\theta}$

\vspace{4pt}
\textit{(Hint: rewrite $\tan\theta$ and $\cot\theta$ in terms of sine and cosine, then combine fractions.)}

\writelines{5}

\vspace{8pt}
\textbf{Part B --- Exact Value.} Find the exact value of $\sin\!\left(\dfrac{5\pi}{12}\right)$.

\textit{(Hint: $\dfrac{5\pi}{12} = \dfrac{\pi}{4} + \dfrac{\pi}{6}$. Apply the sum identity.)}

\writelines{5}

$\sin\!\left(\dfrac{5\pi}{12}\right) = $ \underline{\hspace{4cm}}

\vspace{4pt}
\textbf{Check:} Use a calculator to verify your exact form numerically.

\end{tcolorbox}

\end{document}
